\documentclass[a4paper,11pt]{article}
\usepackage[utf8]{inputenc}
\usepackage[french]{babel}
\usepackage[T1]{fontenc}
\usepackage{amsmath,amssymb,amsthm}
\usepackage{geometry}
\geometry{margin=2cm}
\usepackage{enumitem}
\usepackage{hyperref}
\usepackage{booktabs}
\usepackage{array}
\usepackage{xcolor}
\usepackage{listings}
\usepackage{titlesec}
\usepackage{fancyhdr}
\usepackage{lastpage}
\usepackage{graphicx}
\usepackage{etoolbox}
\AtBeginDocument{\shorthandoff{;:!?}}

% ---- Style des titres ----
\titleformat{\section}{\large\bfseries\color{blue!60!black}}{}{0pt}{}[\vspace{-0.5ex}\rule{\textwidth}{0.4pt}]
\titleformat{\subsection}{\bfseries\color{blue!40!black}}{}{0pt}{}
\titleformat{\subsubsection}{\itshape\color{blue!30!black}}{}{0pt}{}

% ---- Style listings ----
\lstdefinestyle{monstyle}{
  frame=single,
  numbers=left,
  numbersep=5pt,
  breaklines=true,
  basicstyle=\small\ttfamily,
  backgroundcolor=\color{gray!5},
  rulecolor=\color{gray!40},
  columns=fullflexible,
  keepspaces=true,
  showstringspaces=false
}
\lstset{style=monstyle}

% ---- En-têtes / pieds ----
\pagestyle{fancy}
\fancyhf{}
\fancyhead[L]{\small STA211 -- Étude de cas}
\fancyhead[R]{\small Fiche récapitulative}
\fancyfoot[C]{\thepage/\pageref{LastPage}}
\renewcommand{\headrulewidth}{0.4pt}

% ---- Theorem-like environment ----
\theoremstyle{definition}
\newtheorem*{definition}{Définition}
\theoremstyle{remark}
\newtheorem*{remark}{Remarque}
\newtheorem*{important}{À retenir}

\title{\textbf{Fiche récapitulative -- STA211}\\[4pt]
\large Étude de cas : pré-traitement des données German Credit}
\author{}
\date{10 juin 2026}

\begin{document}
\maketitle
\thispagestyle{empty}

\begin{abstract}
Cette fiche présente un cas pratique complet de pré-traitement sur les
données \emph{German Credit} (UCI). Elle illustre le pipeline
exploratoire et les transformations à appliquer avant la modélisation :
importation, analyse univariée et bivariée, transformations de
variables, réduction des données.
\end{abstract}

\vfill
\tableofcontents
\newpage

% ===================================================================
\section{Introduction}
% ===================================================================

L'étude de cas applique les méthodes de pré-traitement vues en cours
aux données \textbf{German Credit} (1\,000 crédits, 21 variables, 20
explicatives + réponse binaire good/bad). L'objectif est de préparer
les données pour une classification supervisée (prédire le statut bon
ou mauvais payeur).

Données~: \url{https://archive.ics.uci.edu/ml/datasets/statlog+(german+credit)}

\subsection{Chargement et librairies}

\begin{lstlisting}[language=R]
library(stargazer)
library(FactoMineR)
library(DescTools)
library(car)
library(lattice)
library(gridExtra)

don <- read.table(
  "https://archive.ics.uci.edu/ml/machine-learning-databases
   /statlog/german/german.data",
  sep = " ", stringsAsFactors = TRUE)
\end{lstlisting}

\subsection{Nettoyage des noms}

Les noms des variables et modalités sont codés (V1, V2, \dots, A11,
A12, \dots). On les renomme d'après le descriptif UCI~:

\begin{lstlisting}[language=R]
colnames(don) <- c("Status", "Duration", "History", "Purpose",
  "Credit.Amount", "Savings account/bonds",
  "Length.of.current.employment", "Instalment.per.cent",
  "Sex.Marital.Status", "Guarantors",
  "Duration.in.Current.address", "Property", "Age.years",
  "Other.installment.plans", "Housing",
  "No.of.Credits.at.this.Bank", "Job", "No.of.dependents",
  "Telephone", "Foreign.Worker", "Creditability")
\end{lstlisting}

On modifie également le type des variables~: \texttt{Duration},
\texttt{Credit.Amount}, \texttt{Age.years} en \texttt{numeric},
et \texttt{Creditability} en \texttt{factor} avec niveaux
\texttt{good}/\texttt{bad}.

% ===================================================================
\section{Aperçu général}
% ===================================================================

\begin{lstlisting}[language=R]
dim(don)               # 1000 x 21
table(sapply(don, class))
# factor:14, integer:4, numeric:3
sum(is.na(don))        # 0 (aucune donnee manquante)
\end{lstlisting}

\begin{important}
  Le jeu de données German Credit est complet (0~ valeur manquante).
  Les étapes de pré-traitement se concentrent donc sur les
  transformations et la réduction, pas sur l'imputation.
\end{important}

% ===================================================================
\section{Analyse univariée}
% ===================================================================

\subsection{Variables quantitatives}

\textbf{Indicateurs statistiques}~: on utilise \texttt{stargazer}
pour un tableau récapitulatif (min, Q1, médiane, moyenne, Q3, max,
écart-type).

\begin{lstlisting}[language=R]
stargazer(don, summary.stat = c("n", "min", "p25", "median",
                                "mean", "p75", "max", "sd"),
          type = "text")
\end{lstlisting}

\textbf{Constats}~:
\begin{itemize}
  \item \texttt{Duration} (4--72 mois, médiane 18)~: asymétrie droite.
  \item \texttt{Credit.Amount} (250--18\,424, médiane 2\,319)~:
    forte asymétrie droite.
  \item \texttt{Age.years} (19--75, médiane 33)~: asymétrie droite
    modérée.
  \item \texttt{Instalment.per.cent}, \texttt{Duration.in.Current.address},
    \texttt{No.of.Credits.at.this.Bank}, \texttt{No.of.dependents}~:
    variables discrètes à peu de modalités.
\end{itemize}

\subsection{Graphiques}

\begin{lstlisting}[language=R]
# Diagrammes en barres (variables discretes)
varbarplot <- c("Instalment.per.cent",
  "Duration.in.Current.address",
  "No.of.Credits.at.this.Bank", "No.of.dependents")
mapply(don[,varbarplot], FUN = function(xx,name)
  {barplot(table(xx), main = name)}, name = varbarplot)

# Histogrammes (variables continues)
varhist <- c("Duration", "Credit.Amount", "Age.years")
mapply(don[,varhist], FUN = function(xx,name)
  {hist(xx, main = name)}, name = varhist)
\end{lstlisting}

Les histogrammes confirment l'asymétrie droite de \texttt{Duration},
\texttt{Credit.Amount} et \texttt{Age.years}. Une transformation
logarithmique ou Box-Cox pourra être envisagée.

\subsection{Boîtes à moustaches}

\begin{lstlisting}[language=R]
library(car)
mapply(don[,varhist], FUN = function(xx,name)
  {Boxplot(xx, main = name, id.n = 2)}, name = varhist)
\end{lstlisting}

Des valeurs apparaissent comme potentiellement aberrantes (notamment
pour \texttt{Credit.Amount}), mais la décision de les traiter pourra
être prise a posteriori si nécessaire.

\subsection{Variables qualitatives}

On détecte les modalités rares via des diagrammes en barres~:

\begin{lstlisting}[language=R]
var.factor <- which(sapply(don, class) == "factor")
mapply(don[,var.factor], FUN = function(xx,name)
  {barplot(table(xx), main = name, horiz = TRUE, las = 2,
           xlim = c(0,1000))},
  name = names(var.factor))
\end{lstlisting}

Exemples de modalités rares~: \texttt{noCredit.allPaid} et
\texttt{thisBank.AllPaid} pour \texttt{History}, \texttt{Guarantor}
pour \texttt{Guarantors}, etc. Ces modalités pourront être regroupées.

% ===================================================================
\section{Analyse bivariée}
% ===================================================================

\subsection{Liens avec la variable réponse}

\textbf{Linéarité (variables quantitatives)}~: on représente la
proportion de bons payeurs en fonction de chaque variable.

\begin{lstlisting}[language=R]
cont.table <- table(don$Duration, don$Creditability)
prof.lignes <- prop.table(cont.table, 1)
plot(as.numeric(rownames(prof.lignes)), prof.lignes[,1],
     pch = 16, xlab = "Duree du credit",
     ylab = "Proportion de bons payeurs")
\end{lstlisting}

La proportion de bons payeurs décroît linéairement avec la durée du
crédit. Pour l'âge, la relation est non-monotone~: une discrétisation
en 3 classes est recommandée.

\textbf{Rapports de corrélation $\eta^2$}~:
\begin{lstlisting}[language=R]
library(BioStatR)
res.eta2 <- sapply(don[,var.numeric], eta2, y = don$Creditability)
barplot(sort(res.eta2), horiz = TRUE, las = 2, xlab = expression(eta^2))
\end{lstlisting}

Les variables les plus discriminantes~: \texttt{Duration} et
\texttt{Credit.Amount}. Les moins discriminantes~:
\texttt{No.of.dependents} et \texttt{Duration.in.Current.address}.

\textbf{V de Cramer (variables qualitatives)}~:
\begin{lstlisting}[language=R]
library(DescTools)
res.cramer <- sapply(don[,var.factor[-which(names(var.factor)
  == "Creditability")]],
  FUN = function(xx) CramerV(table(xx, don$Creditability)))
barplot(sort(res.cramer), horiz = TRUE, las = 2, xlab = "V de Cramer")
\end{lstlisting}

Les plus liées~: \texttt{Status} et \texttt{History}.
Les moins liées~: \texttt{Job} et \texttt{Telephone}.

\begin{important}
  Les variables les moins discriminantes sont candidates à la
  suppression en cas de réduction de dimension.
\end{important}

\subsection{Liens entre variables explicatives}

\textbf{Corrélations (quantitatives)}~:
\begin{lstlisting}[language=R]
matcor <- cor(don[, var.numeric])
PlotCorr(matcor) # visualisation
matcor.spearman <- cor(don[, var.numeric], method = "spearman")
\end{lstlisting}

On compare Pearson et Spearman~: si $|r| \approx |\rho|$, la relation
est linéaire~; si $|\rho| > |r|$, elle est monotone non linéaire.

\textbf{Variables qualitatives}~: on utilise le V de Cramer ou le
$\chi^2$ pour détecter les liaisons fortes entre variables explicatives,
qui pourraient causer de l'instabilité dans les modèles.

% ===================================================================
\section{Analyse multivariée}
% ===================================================================

\subsection{Analyse factorielle}

\textbf{ACM} (pour variables qualitatives) ou \textbf{AFDM} (mixtes)
permettent de visualiser les proximités entre individus et modalités.

\begin{lstlisting}[language=R]
library(FactoMineR)
res.famd <- FAMD(don[, -ncol(don)], graph = FALSE)
plot(res.famd, choix = "var")
plot(res.famd, choix = "ind")
\end{lstlisting}

L'AFDM permet également de détecter des individus atypiques
(contribution élevée à une dimension) et des liaisons entre
variables qualitatives et quantitatives.

\subsection{Classification}

\begin{lstlisting}[language=R]
# CAH sur les composantes de l'AFDM
res.hcpc <- HCPC(res.famd, nb.clust = -1, graph = FALSE)
plot(res.hcpc)
\end{lstlisting}

La classification permet d'identifier des profils types de clients
(segments homogènes) et de visualiser les caractéristiques de
chaque cluster via la fonction \texttt{catdes}.

\subsection{Description des partitions}

\begin{lstlisting}[language=R]
catdes(don, num.var = ncol(don))
\end{lstlisting}

La fonction \texttt{catdes} décrit chaque modalité de la variable
réponse (\texttt{good}/\texttt{bad}) à partir~:
\begin{itemize}
  \item Des variables quantitatives (test de Student, moyenne dans
    chaque groupe).
  \item Des modalités des variables qualitatives (test du $\chi^2$,
    proportion dans chaque groupe).
\end{itemize}

% ===================================================================
\section{Transformations}
% ===================================================================

\subsection{Variables quantitatives}

\textbf{Discrétisation}~: on coupe l'âge en 3 classes pour gérer
la non-linéarité observée~:

\begin{lstlisting}[language=R]
Age.cut <- cut(don$Age.years,
  breaks = quantile(don$Age.years, probs = c(0, 1/3, 2/3, 1)),
  include.lowest = TRUE)
table(Age.cut)
\end{lstlisting}

\textbf{Box-Cox}~: pour linéariser et symétriser les distributions.

\begin{lstlisting}[language=R]
library(MASS)
boxcox(lm(Credit.Amount ~ 1, data = don))
\end{lstlisting}

La transformation logarithmique ($\lambda \approx 0$) est souvent
appropriée pour \texttt{Credit.Amount} et \texttt{Duration}.

\subsection{Variables qualitatives}

\textbf{Regroupement de modalités}~: on fusionne les modalités rares
qui ont un comportement similaire vis-à-vis de la réponse.

Exemple~: \texttt{noCredit.allPaid} et \texttt{thisBank.AllPaid} de
\texttt{History} ont des proportions de bons payeurs proches~:

\begin{lstlisting}[language=R]
tmp <- table(don$History, don$Creditability)
prop.table(tmp, 1)  # comparer les profils
levels(don$History)[1:2] <- "AllPaid"
\end{lstlisting}

\subsection{Standardisation}

Pour les méthodes sensibles aux échelles (ACP, régression, $k$-NN),
on centre-réduit les variables quantitatives~:

\begin{lstlisting}[language=R]
don.scaled <- as.data.frame(scale(don[, var.numeric]))
\end{lstlisting}

% ===================================================================
\section{Réduction des données}
% ===================================================================

\subsection{Réduction en lignes}

On peut supprimer les individus atypiques identifiés par~:
\begin{itemize}
  \item Distance de Mahalanobis robuste (MCD).
  \item Isolation Forest (peu sensible à la dimension).
\end{itemize}

\begin{lstlisting}[language=R]
library(robustbase)
mcd <- covMcd(don[, var.numeric])
dist <- mahalanobis(don[, var.numeric],
                    center = mcd$center, cov = mcd$cov)
seuil <- qchisq(0.999, df = length(var.numeric))
outliers <- which(dist > seuil)
\end{lstlisting}

\subsection{Réduction en colonnes}

\textbf{Sélection de variables}~:
\begin{itemize}
  \item Éliminer les variables non discriminantes (faible $\eta^2$,
    faible V de Cramer).
  \item Éliminer les variables redondantes (forte corrélation entre
    elles).
  \item Utiliser le Lasso pour une sélection automatique~:
\end{itemize}

\begin{lstlisting}[language=R]
library(glmnet)
X <- model.matrix(Creditability ~ ., data = don)[, -1]
y <- don$Creditability
cv.lasso <- cv.glmnet(X, y, alpha = 1, family = "binomial")
coef(cv.lasso, s = "lambda.1se")
\end{lstlisting}

\textbf{Extraction de caractéristiques}~:
\begin{itemize}
  \item ACP (variables numériques), ACM (qualitatives), AFDM (mixtes).
  \item On conserve les premières composantes qui résument
    l'essentiel de l'information.
\end{itemize}

\begin{lstlisting}[language=R]
res.pca <- PCA(don[, var.numeric], scale.unit = TRUE, graph = FALSE)
barplot(res.pca$eig[, 2], main = "Variance expliquee")
# Souvent 2-3 composantes suffisent
\end{lstlisting}

% ===================================================================
\section{Résumé du pipeline}
% ===================================================================

\begin{enumerate}
  \item \textbf{Import}~: charger, renommer, typer.
  \item \textbf{Aperçu}~: dimensions, types, valeurs manquantes.
  \item \textbf{Analyse univariée}~: distributions, atypiques,
    modalités rares.
  \item \textbf{Analyse bivariée}~: lien avec $Y$ (linéarité,
    $\eta^2$, V de Cramer), lien entre $X$.
  \item \textbf{Analyse multivariée}~: AFDM, CAH, catdes.
  \item \textbf{Transformations}~: Box-Cox, discrétisation,
    regroupement, standardisation.
  \item \textbf{Réduction}~: suppression d'atypiques, sélection de
    variables (Lasso), extraction (ACP/AFDM).
  \item \textbf{Modélisation}~: appliquer la méthode choisie sur
    les données prétraitées.
\end{enumerate}

% ===================================================================
\section*{Questions clés (QCM)}
% ===================================================================

\addcontentsline{toc}{section}{Questions clés (QCM)}

\begin{enumerate}

  \item Dans l'étude de cas German Credit, combien de variables
    le jeu de données contient-il~?
    \begin{enumerate}
      \item 1000
      \item 21
      \item 20
      \item 14
    \end{enumerate}
    \textbf{Réponse~: b (21 variables dont 20 explicatives + réponse).}

  \item Le jeu de données German Credit contient-il des valeurs
    manquantes~?
    \begin{enumerate}
      \item Oui, beaucoup
      \item Non, il est complet
      \item Oui, 5\%
      \item On ne peut pas savoir
    \end{enumerate}
    \textbf{Réponse~: b (sum(is.na(don)) = 0).}

  \item Quelle variable quantitative est la plus discriminante
    (liée à la réponse) selon le $\eta^2$~?
    \begin{enumerate}
      \item Age.years
      \item Duration
      \item Instalment.per.cent
      \item No.of.dependents
    \end{enumerate}
    \textbf{Réponse~: b (Duration).}

  \item Quelle variable qualitative est la plus discriminante selon
    le V de Cramer~?
    \begin{enumerate}
      \item Telephone
      \item Job
      \item Status
      \item Housing
    \end{enumerate}
    \textbf{Réponse~: c (Status).}

  \item La relation entre Age et la proportion de bons payeurs est~:
    \begin{enumerate}
      \item Linéaire croissante
      \item Linéaire décroissante
      \item Non monotone
      \item Il n'y a pas de relation
    \end{enumerate}
    \textbf{Réponse~: c (non monotone, d'où la discrétisation en 3 classes).}

  \item Quel package R permet de calculer le $\eta^2$~?
    \begin{enumerate}
      \item stargazer
      \item DescTools
      \item BioStatR
      \item FactoMineR
    \end{enumerate}
    \textbf{Réponse~: c (BioStatR, fonction eta2).}

  \item La fonction \texttt{catdes} du package FactoMineR permet de~:
    \begin{enumerate}
      \item Visualiser les corrélations
      \item Décrire une partition par les variables
      \item Effectuer une ACP
      \item Imputer les valeurs manquantes
    \end{enumerate}
    \textbf{Réponse~: b (description d'une partition).}

  \item Pour gérer la non-linéarité de l'âge, on recommande~:
    \begin{enumerate}
      \item Une transformation logarithmique
      \item Une discrétisation en 3 classes
      \item Une standardisation
      \item La suppression de la variable
    \end{enumerate}
    \textbf{Réponse~: b (discrétisation en 3 classes par quantiles).}

  \item Le Lasso (régression pénalisée L1) permet~:
    \begin{enumerate}
      \item D'imputer les données manquantes
      \item De sélectionner automatiquement les variables
      \item De visualiser les clusters
      \item De transformer les variables
    \end{enumerate}
    \textbf{Réponse~: b (sélection de variables via pénalité L1).}

  \item Quel est l'ordre recommandé du pipeline de pré-traitement~?
    \begin{enumerate}
      \item Modélisation $\rightarrow$ transformation $\rightarrow$ analyse
      \item Analyse $\rightarrow$ transformation $\rightarrow$ réduction
        $\rightarrow$ modélisation
      \item Réduction $\rightarrow$ modélisation $\rightarrow$ analyse
      \item Transformation $\rightarrow$ modélisation $\rightarrow$ analyse
    \end{enumerate}
    \textbf{Réponse~: b (analyse exploratoire $\rightarrow$ transformations
    $\rightarrow$ réduction $\rightarrow$ modélisation).}

\end{enumerate}

% ===================================================================
\begin{thebibliography}{9}
% ===================================================================

\bibitem{german}
  UCI Machine Learning Repository.
  \newblock \emph{German Credit Data}.
  \newblock \url{https://archive.ics.uci.edu/ml/datasets/statlog+(german+credit)}

\bibitem{audigier2026}
  V.~Audigier, N.~Niang.
  \newblock STA211~: pré-traitement, étude de cas.
  \newblock Cours CNAM, 2026.

\bibitem{tuffery2007}
  S.~Tufféry.
  \newblock \emph{Data Mining et statistique décisionnelle}.
  \newblock Technip, 2007.

\bibitem{lejeune2010}
  M.~Lejeune.
  \newblock \emph{Statistique~: la théorie et ses applications}.
  \newblock Springer, 2010.

\bibitem{lebart2006}
  L.~Lebart, A.~Morineau, M.~Piron.
  \newblock \emph{Statistique exploratoire multidimensionnelle}.
  \newblock Dunod, 2006.

\end{thebibliography}

\end{document}
